\section{Limits and colimits}
\label{section-limits}

\noindent
Let $\mathcal{C}$ be a category. A {\it diagram} in $\mathcal{C}$ is
simply a functor $M : \mathcal{I} \to \mathcal{C}$. We say that
$\mathcal{I}$ is the {\it index category} or that $M$ is an
$\mathcal{I}$-diagram. We will use the notation $M_i$ to denote the
image of the object
$i$ of $\mathcal{I}$. Hence for $\phi : i \to i'$ a morphism
in $\mathcal{I}$ we have $M(\phi) : M_i \to M_{i'}$.

\begin{definition}
\label{definition-limit}
A {\it limit} of the $\mathcal{I}$-diagram $M$ in the category
$\mathcal{C}$ is given by an object $\lim_\mathcal{I} M$ in $\mathcal{C}$
together with morphisms $p_i : \lim_\mathcal{I} M \to M_i$ such that
\begin{enumerate}
\item for $\phi : i \to i'$ a morphism
in $\mathcal{I}$ we have $p_{i'} = M(\phi) \circ p_i$, and
\item for any object $W$ in $\mathcal{C}$ and any family of
morphisms $q_i : W \to M_i$ (indexed by $i \in \Ob(\mathcal{I})$)
such that for all $\phi : i \to i'$
in $\mathcal{I}$ we have $q_{i'} = M(\phi) \circ q_i$ there
exists a unique morphism $q : W \to \lim_\mathcal{I} M$ such that
$q_i = p_i \circ q$ for every object $i$ of $\mathcal{I}$.
\end{enumerate}
\end{definition}

\noindent
Limits $(\lim_\mathcal{I} M, (p_i)_{i\in \Ob(\mathcal{I})})$ are
(if they exist)
unique up to unique isomorphism by the uniqueness requirement
in the definition. Products of pairs, fibre products, and equalizers are
examples of limits. The limit over the empty diagram is a final object
of $\mathcal{C}$.
In the category of sets all limits exist.
The dual notion is that of colimits.

\begin{definition}
\label{definition-colimit}
A {\it colimit} of the $\mathcal{I}$-diagram $M$ in the category
$\mathcal{C}$ is given by an object $\colim_\mathcal{I} M$ in $\mathcal{C}$
together with morphisms $s_i : M_i \to \colim_\mathcal{I} M$ such that
\begin{enumerate}
\item for $\phi : i \to i'$ a morphism
in $\mathcal{I}$ we have $s_i = s_{i'} \circ M(\phi)$, and
\item for any object $W$ in $\mathcal{C}$ and any family of
morphisms $t_i : M_i \to W$ (indexed by $i \in \Ob(\mathcal{I})$)
such that for all $\phi : i \to i'$
in $\mathcal{I}$ we have $t_i = t_{i'} \circ M(\phi)$ there
exists a unique morphism $t : \colim_\mathcal{I} M \to W$ such that
$t_i = t \circ s_i$ for every object $i$ of $\mathcal{I}$.
\end{enumerate}
\end{definition}

\noindent
Colimits $(\colim_\mathcal{I} M, (s_i)_{i\in \Ob(\mathcal{I})})$ are
(if they exist) unique up to unique isomorphism by the uniqueness requirement
in the definition. Coproducts of pairs, pushouts, and coequalizers are
examples of colimits. The colimit over an empty diagram is an initial object
of $\mathcal{C}$. In the category of sets all colimits exist.

\begin{remark}
\label{remark-diagram-small}
The index category of a (co)limit will never be allowed to have
a proper class of objects. In this project it means that
it cannot be one of the categories listed in
Remark \ref{remark-big-categories}
\end{remark}

\begin{remark}
\label{remark-limit-colim}
We often write $\lim_i M_i$, $\colim_i M_i$,
$\lim_{i\in \mathcal{I}} M_i$, or $\colim_{i\in \mathcal{I}} M_i$
instead of the versions indexed by $\mathcal{I}$.
Using this notation, and using the description of
limits and colimits of sets in Section \ref{section-limit-sets}
below, we can say the following.
Let $M : \mathcal{I} \to \mathcal{C}$ be a diagram.
\begin{enumerate}
\item The object $\lim_i M_i$ if it exists satisfies the following property
$$
\Mor_\mathcal{C}(W, \lim_i M_i)
=
\lim_i \Mor_\mathcal{C}(W, M_i)
$$
where the limit on the right takes place in the category of sets.
\item The object $\colim_i M_i$ if it
exists satisfies the following property
$$
\Mor_\mathcal{C}(\colim_i M_i, W)
=
\lim_{i \in \mathcal{I}^{opp}} \Mor_\mathcal{C}(M_i, W)
$$
where on the right we have the limit over the opposite category
with value in the category of sets.
\end{enumerate}
By the Yoneda lemma (and its dual) this formula completely determines the
limit, respectively the colimit.
\end{remark}

\begin{remark}
\label{remark-cones-and-cocones}
Let $M : \mathcal{I} \to \mathcal{C}$ be a diagram. In this setting a
{\it cone} for $M$ is given by an object $W$ and a family of morphisms
$q_i : W \to M_i$, $i \in \Ob(\mathcal{I})$ such that for all morphisms
$\phi : i \to i'$ of $\mathcal{I}$ the diagram
$$
\xymatrix{
& W \ar[dl]_{q_i} \ar[dr]^{q_{i'}} \\
M_i \ar[rr]^{M(\phi)} & & M_{i'}
}
$$
is commutative. The collection of cones forms a category with an obvious
notion of morphisms. Clearly, the limit of $M$, if it exists, is a final
object in the category of cones. Dually, a {\it cocone} for $M$ is given
by an object $W$ and a family of morphisms $t_i : M_i \to W$ such that for
all morphisms $\phi : i \to i'$ in $\mathcal{I}$ the diagram
$$
\xymatrix{
M_i \ar[rr]^{M(\phi)} \ar[dr]_{t_i} & & M_{i'} \ar[dl]^{t_{i'}} \\
& W
}
$$
commutes. The collection of cocones forms a category with an obvious notion
of morphisms. Similarly to the above the colimit of $M$ exists
if and only if the category of cocones has an initial object.
\end{remark}

\noindent
As an application of the notions of limits and colimits
we define products and coproducts.

\begin{definition}
\label{definition-product}
Suppose that $I$ is a set, and suppose given for every $i \in I$ an
object $M_i$ of the category $\mathcal{C}$. A {\it product}
$\prod_{i\in I} M_i$ is by definition $\lim_\mathcal{I} M$
(if it exists)
where $\mathcal{I}$ is the category having only identities as
morphisms and having the elements of $I$ as objects.
\end{definition}

\noindent
An important special case is where $I = \emptyset$ in which case the
product is a final object of the category.
The morphisms $p_i : \prod M_i \to M_i$ are called the
{\it projection morphisms}.

\begin{definition}
\label{definition-coproduct}
Suppose that $I$ is a set, and suppose given for every $i \in I$ an
object $M_i$ of the category $\mathcal{C}$. A {\it coproduct}
$\coprod_{i\in I} M_i$ is by definition $\colim_\mathcal{I} M$
(if it exists) where $\mathcal{I}$ is the category having only
identities as morphisms and having the elements of $I$ as objects.
\end{definition}

\noindent
An important special case is where $I = \emptyset$ in which case the
coproduct is an initial object of the category.
Note that the coproduct comes equipped with morphisms
$M_i \to \coprod M_i$. These are sometimes called the
{\it coprojections}.

\begin{lemma}
\label{lemma-functorial-colimit}
Suppose that $M : \mathcal{I} \to \mathcal{C}$,
and $N : \mathcal{J} \to \mathcal{C}$ are diagrams
whose colimits exist. Suppose
$H : \mathcal{I} \to \mathcal{J}$ is
a functor, and suppose $t : M \to N \circ H$
is a transformation of functors.
Then there is a unique morphism
$$
\theta :
\colim_\mathcal{I} M
\longrightarrow
\colim_\mathcal{J} N
$$
such that all the diagrams
$$
\xymatrix{
M_i \ar[d]_{t_i} \ar[r]
&
\colim_\mathcal{I} M \ar[d]^{\theta}
\\
N_{H(i)} \ar[r]
&
\colim_\mathcal{J} N
}
$$
commute.
\end{lemma}

\begin{proof}
Omitted.
\end{proof}

\begin{lemma}
\label{lemma-functorial-limit}
Suppose that $M : \mathcal{I} \to \mathcal{C}$,
and $N : \mathcal{J} \to \mathcal{C}$ are diagrams
whose limits exist. Suppose $H : \mathcal{I} \to \mathcal{J}$ is
a functor, and suppose $t : N \circ H \to M$
is a transformation of functors.
Then there is a unique morphism
$$
\theta :
\lim_\mathcal{J} N
\longrightarrow
\lim_\mathcal{I} M
$$
such that all the diagrams
$$
\xymatrix{
\lim_\mathcal{J} N \ar[d]^{\theta} \ar[r]
&
N_{H(i)} \ar[d]_{t_i}
\\
\lim_\mathcal{I} M \ar[r]
&
M_i
}
$$
commute.
\end{lemma}

\begin{proof}
Omitted.
\end{proof}


\begin{lemma}
\label{lemma-colimits-commute}
Let $\mathcal{I}$, $\mathcal{J}$ be index categories. Let
$M : \mathcal{I} \times \mathcal{J} \to \mathcal{C}$ be a functor.
Assume that $M_{i, \infty} = \colim_j M_{i,j}$ 
exists for all $i$. Then the resulting functor
$M_{-, \infty} : \mathcal{I} \to \mathcal{C}$
has a colimit if and only if $M$ does, and then the colimits coincide.
In particular, we have
$$
\colim_i \colim_j M_{i, j}
=
\colim_{i, j} M_{i, j}
=
\colim_j \colim_i M_{i, j}
$$
provided all the indicated colimits exist. Similar for limits.
\end{lemma}

\begin{proof}
Omitted.
\end{proof}

\begin{lemma}
\label{lemma-limits-products-equalizers}
Let $M : \mathcal{I} \to \mathcal{C}$ be a diagram.
Write $I = \Ob(\mathcal{I})$ and $A = \text{Arrows}(\mathcal{I})$.
Denote $s, t : A \to I$ the source and target maps.
Suppose that $\prod_{i \in I} M_i$ and $\prod_{a \in A} M_{t(a)}$
exist. Suppose that the equalizer of
$$
\xymatrix{
\prod_{i \in I} M_i
\ar@<1ex>[r]^\phi \ar@<-1ex>[r]_\psi
&
\prod_{a \in A} M_{t(a)}
}
$$
exists, where the morphisms are determined by their components
as follows: $p_a \circ \psi = M(a) \circ p_{s(a)}$
and $p_a \circ \phi = p_{t(a)}$. Then this equalizer is the
limit of the diagram.
\end{lemma}

\begin{proof}
Omitted.
\end{proof}


\begin{lemma}
\label{lemma-colimits-coproducts-coequalizers}
\begin{slogan}
If all coproducts and coequalizers exist, all colimits exist.
\end{slogan}
Let $M : \mathcal{I} \to \mathcal{C}$ be a diagram.
Write $I = \Ob(\mathcal{I})$ and $A = \text{Arrows}(\mathcal{I})$.
Denote $s, t : A \to I$ the source and target maps.
Suppose that $\coprod_{i \in I} M_i$ and $\coprod_{a \in A} M_{s(a)}$
exist. Suppose that the coequalizer of
$$
\xymatrix{
\coprod_{a \in A} M_{s(a)}
\ar@<1ex>[r]^\phi \ar@<-1ex>[r]_\psi
&
\coprod_{i \in I} M_i
}
$$
exists, where the morphisms are determined by their components
as follows: The component $M_{s(a)}$ maps via $\psi$
to the component $M_{t(a)}$ via the morphism $M(a)$.
The component $M_{s(a)}$ maps via $\phi$ to the component
$M_{s(a)}$ by the identity morphism. Then this coequalizer is the
colimit of the diagram.
\end{lemma}

\begin{proof}
Omitted.
\end{proof}
